% ============================================================
% Model Evaluation: Strengths and Weaknesses (Concise Version)
% ============================================================
\section{Strengths and Limitations}
\label{sec:evaluation}

% ============================================================
% Model Evaluation: Strengths and Weaknesses (Compressed Layout)
% ============================================================



% 开启小字号模式，并缩减段后间距
\begingroup

\setlength{\parskip}{0pt} 

We assess the validity of our framework by evaluating its operational identifiability, structural consistency, and interpretability under partial observability.

%\vspace{-0.8em} % 强制缩减标题与列表的距离
\noindent\textbf{Strengths.}
\begin{enumerate}
    % 强制压缩列表内部间距
    \setlength{\itemsep}{0pt}
    \setlength{\parskip}{0pt}
    \setlength{\parsep}{0pt}
    
    \item \textbf{Unified Uncertainty Propagation.} 
    By propagating the full posterior distribution ($M=2000$ draws) rather than point estimates, our framework ensures operational identifiability and avoids the false precision of deterministic models ($\mathrm{PPC} \approx 0.905$).

    %\item \textbf{Structural Isolation via Counterfactuals.} 
    %Our within-week design isolates the aggregation rule as the sole independent variable, confirming that observed outcome divergences are inherent systemic properties rather than random anomalies.

    \item \textbf{Robust Dual-Signal Attribution.} 
    The hazard model effectively decouples meritocratic (judge) from democratic (fan) signals, with parametric findings validated by the consistent feature rankings of nonparametric random forests.

    \item \textbf{Stress-Tested Mechanism Design.} 
    The proposed PAGS architecture is validated against historical "worst-case" controversies, demonstrating quantifiable fairness improvements without sacrificing commercial audience agency.
\end{enumerate}

%\vspace{-0.8em} % 强制缩减段落间距
\noindent\textbf{Limitations.}
\begin{enumerate}
    % 强制压缩列表内部间距
    \setlength{\itemsep}{0pt}
    \setlength{\parskip}{0pt}
    \setlength{\parsep}{0pt}

    \item \textbf{Dependence on Latent Inference.} 
    All findings remain conditional on the inverse reconstruction; weeks with sparse data or weak correlations yield diffuse posteriors that inherently widen downstream uncertainty bands.

    \item \textbf{Static Behavioral Assumption.} 
    The counterfactuals assume a \textit{ceteris paribus} environment, isolating mechanical effects while neglecting potential long-term game-theoretic adaptations by judges or contestants.

    \item \textbf{Parametric Sensitivity.} 
    The proposed system relies on tuned damping parameters, which require continuous auditing to ensure robustness against rare, extreme voting distributions in future seasons.
\end{enumerate}
\vspace{1ex}
\noindent\textbf{Conclusion.}  

Overall, this work establishes an uncertainty-aware audit framework for DWTS-style hybrid decision rules under unobserved voting. With regime-consistent vote reconstruction (PPC $=90.5\%$), quantified rule-induced volatility ($26\%$ of historical weeks), and cross-model driver-invariance evidence, we provide a systematic diagnostic of the conditions and origins of legitimacy risks. PAGS maintains a near-neutral performance baseline ($P \approx 0.5$) to preserve audience agency while enforcing a transparent circuit-breaker that corrects extreme popularity--merit divergences (e.g., the Bobby Bones corner case), and the producer-facing memo translates these audit findings into deployable guidance.


\endgroup